\section{The exact remaining asymptotic target}
\label{ct:remaining}

\subsection{Minimality already removes the rounding problem}

Fix \(c>2\), an ambient \(M\), and an interval \(I\). If an
instance fails \(\tau(A,I)\le c|A|\), replace \(A\) by a failing
subset of smallest cardinality, set \(m=|A|\), and write
\[
 K=\lfloor c(m-1)\rfloor+1.
\]
Every deletion has an integer cover of size at most \(K-1\).
Adding two unused positions would give a cover within the budget
\(\lfloor cm\rfloor\), so the forced-bank arguments apply.
For completeness, their short proofs are retained here.

For every prime occurrence \(q\mid a\), all interval multiples
of \(a/q\) belong to every deletion cover of \(A\setminus\{a\}\).
Otherwise a fresh multiple supplies \(a/q\), and at most one extra
unused position repairs the missing \(q\)-unit. Full supply
\eqref{ct:supply} guarantees that repair. Since there are at least
\(q\) such multiples, all source primes satisfy
\begin{equation}
 p<K<cm.
 \label{ct:prime-cap}
\end{equation}

If \(pq\mid a\), counting occurrences and allowing \(p=q\),
and \(pq\ge K\), then
\begin{equation}
 D_p>v_p\left(\prod_{\substack{b\in I\\a/q\mid b}}b\right)
 \ge qv_p(a/q).
 \label{ct:pair-bank}
\end{equation}
If the first strict inequality failed, the forced bank would
already pay the full \(p\)-demand. There are at least \(pq\ge K\)
targets divisible by \(a/(pq)\); one is fresh relative to a cover
of size \(K-1\). Add it, then repair at most one \(q\)-unit.
This again contradicts minimal failure.

Define the demand core and its complement by
\begin{equation}
 E=\{p\in P:p\le\sqrt K\}\cup\{p\in P:D_p>p\},
 \qquad R=P\setminus E.
 \label{ct:demand-core}
\end{equation}
Every source has at most one \emph{occurrence} from \(R\), and
that prime is its unique largest prime. Indeed, if \(r\in R\)
and another occurrence \(q\ge r\) were in the same source, then
\(rq>K\), while \(D_r\le r\le qv_r(a/q)\), contrary to
\eqref{ct:pair-bank}. In particular,
\begin{equation}
 R\subset(\sqrt K,K),\qquad
 |P|\le\pi(cm)=O_c(m/\log m).
 \label{ct:round-small}
\end{equation}

The quadratic gate forces \(M<16m^2\). For
\(H=\{p:D_p>p\}\), the logarithmic demand budget is
\[
 \sum_{p\in H}p\log p
 <\sum_pD_p\log p
 =\sum_{a\in A}\log a
 \le m\log M.
\]
Chebyshev's bound gives \(p_j\gg j\log(j+1)\) for the
\(j\)-th prime. Thus the left side, for \(h=|H|\), is at least
a constant times \(h^2\log^2(h+1)\). Separating
\(h<m^{1/4}\) from its complement gives
\begin{equation}
 |E|=O_c\bigl(\sqrt{m/\log m}\bigr).
 \label{ct:core-size}
\end{equation}
This reproduces the sharper structural bound in \cite{Synthesis}
without an additive-budget parameter. Minimality is always taken
at fixed \(M\); no interval is shortened after a deletion.

\subsection{A sufficient inequality for full asymptotic closure}

Fix \(\delta>0\) and take \(c=2+3\delta\) in the preceding
minimality argument. The remaining analytic goal is the following
uniform estimate for the just-described minimal-obstruction class:
\begin{equation}
 \boxed{\quad
 \sup_{0\le w_p\le1}\Delta_{2+\delta}(A,I,w)=o(m).
 \quad}
 \label{ct:open-hinge}
\end{equation}
The supremum is over \emph{all} clipped core weights, with no
bound on how many exceed a small threshold. The estimate must be
uniform in \(M<16m^2\), the interval shift, the distinct sources,
and their prime arithmetic.

To see that this goal suffices, suppose arbitrarily large instances
failed the budget \((2+3\delta)n\), and pass at fixed ambient
length to a minimal failing subset of size \(m\). The quadratic
gate gives \(m>\sqrt M/4\), so \(m\to\infty\). By
\eqref{ct:open-hinge} and duality,
\(L(A,I)\le(2+\delta)m+o(m)\). By
\eqref{ct:round-small} and \eqref{ct:round}, the same bound holds
for \(\tau(A,I)\), contradicting the chosen budget for large
\(m\). Since \(\delta\) is arbitrary, the unrestricted
\((2+o(1))n\) assertion would follow.

The new controlled-core theorem bounds the supremum restricted to
vectors having at most
\[
 k\le \frac{(1-\gamma)\log\log m}{\log D_\delta},
 \qquad 0<\gamma<1,
\]
core weights exceeding \(c_\delta/(k+1)\) by \(o(m)\). Arbitrarily many
substantial weights in \(R\) are allowed. The unrestricted demand
core in \eqref{ct:core-size} can be much larger, so the result
does not establish \eqref{ct:open-hinge} for every vector.

The separate moment theorem excludes dual vectors with mean
\(\sum_{p\le M}w_p/(p-1)\ge241/136+o(1)\).
The separated-support theorem excludes its stated source classes.
Combining these exclusions does not show that every remaining
instance belongs to one of them.

The direct construction and its strengthening in
Corollary~\ref{cor:sc-chebyshev-core} give a further source-level
restriction, independent of the weights. For the sufficiently
small absolute \(c_0>0\) of that corollary, the demand core of a
sufficiently large minimal asymptotic obstruction must satisfy
\[
 |E|>c_0\frac{(\log M)^2}{\log\log M}.
\]
Indeed, the one-outside-prime property already holds by
\eqref{ct:demand-core}, so a smaller core would place that source
set inside the new integer theorem. Consequently any hypothetical
obstruction sequence has
\[
 \liminf\frac{\log|E|}{\log\log M}\ge2,
 \qquad |E|=O_c\bigl(\sqrt{m/\log m}\bigr).
\]
There is still a large gap between these necessary bounds. This is
not a proof that such an obstruction exists.

\subsection{What has been closed and what still has to change}

The following steps are now established within this proof programme:

\begin{enumerate}
\item The diffuse core and the one-prime remainder share a single
target budget. No separate core cover is added to a remainder cover.
\item All raw prime-power tails are retained. Capping is an exact
intermediate device, followed by a justified simultaneous restoration.
\item A controlled substantial core incurs a finite upper-set error,
independent of \(M\). Each remainder prime pays that error once,
regardless of its exponent range.
\item Fixed-ambient minimality reduces integer rounding to \(o(m)\).
The raw-to-integer passage is no longer an unproved final step.
\item A direct common-core construction handles the core size
bound \(c_0(\log M)^2/\log\log M\), while retaining arbitrary
outside prime powers. Nested matching supplies all outside demands;
exact core-part bins and a sparse exceptional repair supply the
shared core demands.
\item In the separated-support regime, arbitrary pure-core sources
are absorbed, positive certificate overlap is bounded, and the
old positive lower bound on the fixed large-prime exponent is removed.
\end{enumerate}

For the unrestricted asymptotic target, the next proof must control
the joint contribution of many substantial core coordinates more
economically than the variation bound
\(B_{\delta,k}(1+|R|)\). It is sufficient to prove
\eqref{ct:open-hinge}; it is not necessary to bound each upper-set
discrepancy separately. Any such proof must preserve the shared
target budget and the dependence between the core, the remainders,
and the source demands.

For the exact target \(2n\), even a proof of
\eqref{ct:open-hinge} would leave an additional task: removing the
\(o(n)\) slack and the rounding loss, or constructing an exact
integer cover directly. The present note therefore distinguishes
the exact and asymptotic endpoints throughout.

\subsection{Two shortcuts that do not establish the missing step}

The upper set of smooth parts cannot always be counted in a prefix
as though that prefix minimized its count among intervals of equal
length. The set of integers divisible by \(2\) or \(3\) contains
two members of \(\{1,2,3\}\), but only one member of
\(\{5,6,7\}\). This disproves that intermediate count claim;
it is not a counterexample to problem 708. The finite-difference
error in the controlled-core proof is what repairs this step.

Also, the statement ``for every fixed \(\alpha>0\)'' in the
separated-support theorem cannot be used with an arbitrary
instance-dependent \(\alpha\to0\). The shrinking-exponent example
is valid because its quantitative error conditions are checked
separately. Primes between its stated bounds \(Y(M)\) and \(Z(M)\)
remain outside that source class.

The archive contains the complete proof notes, their mathematical
audit records, and the inherited inputs. No finite search is used
as evidence of a universal conclusion in this continuation.
